% ============================================================================
%  ESTADÍSTICA DESCRIPTIVA — Sesión 2
%  Semana 2: 17 al 22 de agosto de 2026   (lunes 17: festivo)
%  1.2 Población, muestra, parámetros y estadísticos
%  1.3 Tipos de variables · 1.4 Escalas de medición
% ============================================================================
\documentclass[aspectratio=169,10pt]{beamer}

\usepackage[utf8]{inputenc}
\IfFileExists{spanish.ldf}{\usepackage[spanish,es-noshorthands]{babel}}{}
\usepackage{amsmath,amssymb}
\usepackage{booktabs}
\usepackage{tikz}
\usepackage{listings}

\usetheme{Uniminuto}

\uniasignatura{Estadística descriptiva}
\uniprofesor{Andrés Felipe Puerta Velez}
\unicorreo{andres.puerta-v@uniminuto.edu.co}
\uninrc{CBASBL031 - CBASBL141}
\unisemana{Sesión 2 — Semana 2 · 17 al 22 de agosto de 2026}


\newcommand{\fx}[1]{{\ttfamily\color{uniAzul}#1}}
\newcommand{\listatiny}{\setbeamerfont{itemize/enumerate body}{size=\tiny}\tiny}

% ------------------------------------------------------------- CÓDIGO R ----
\lstdefinestyle{Rmini}{
  language=R,
  basicstyle=\ttfamily\tiny,
  keywordstyle=\color{uniAzul}\bfseries,
  commentstyle=\color{uniGris}\itshape,
  stringstyle=\color{uniTitulo},
  showstringspaces=false, numbers=none, breaklines=true,
  columns=fullflexible,
  backgroundcolor=\color{uniCremaClaro},
  frame=leftline, framerule=1.2pt, rulecolor=\color{uniNaranja},
  xleftmargin=0.35cm, framexleftmargin=0.35cm,
  framexrightmargin=0.15cm, framextopmargin=2pt, framexbottommargin=2pt,
  aboveskip=0.12cm, belowskip=0.12cm
}
\lstdefinestyle{Rout}{
  basicstyle=\ttfamily\tiny\color{black},
  showstringspaces=false, numbers=none, breaklines=true,
  columns=fullflexible,
  backgroundcolor=\color{uniGris!10},
  frame=leftline, framerule=1.2pt, rulecolor=\color{uniGris},
  xleftmargin=0.35cm, framexleftmargin=0.35cm,
  framexrightmargin=0.15cm, framextopmargin=2pt, framexbottommargin=2pt,
  aboveskip=0.12cm, belowskip=0.12cm
}

\begin{document}

\portada

\begin{frame}{Orden del día}
  \begin{enumerate}
    \item Llamado a lista y revisión de los ejercicios propuestos
    \item Repaso: dónde quedamos
    \item \azul{Tema 1.2: parámetros y estadísticos}
    \item \azul{Tema 1.3: tipos de variables} y los números que no son cantidades
    \item \azul{Tema 1.4: escalas de medición} y qué permite cada una
    \item Taller de clasificación, errores frecuentes y cierre
  \end{enumerate}

  \vspace{0.25cm}
  {\footnotesize\color{uniGris} Sesión de 3 horas. Lo de hoy decide
  \destacar{qué tabla, qué gráfico y qué medida} se podrán usar en el resto
  del curso.}
\end{frame}

\begin{frame}{Repaso: dónde quedamos}
  En la sesión 1 quedó claro qué son población, muestra, unidad de observación,
  variable y dato.

  \vspace{0.25cm}
  \begin{alertblock}{Las preguntas de hoy}
    \footnotesize
    Si calculo un promedio, ¿ese número describe a la población o a la muestra?
    Y sobre todo: \destacar{¿todas las variables se pueden tratar igual?} La
    respuesta a la segunda determina qué tabla, qué gráfico y qué medida se
    puede usar en el resto del curso.
  \end{alertblock}
\end{frame}

\seccion{1.2 Parámetros y estadísticos}

\begin{frame}{La distinción que hay que memorizar}
  \begin{columns}[T,onlytextwidth]
    \begin{column}{0.48\textwidth}
      \begin{block}{Parámetro}
        \footnotesize
        Número que describe a la \azul{población}. Es fijo y casi siempre
        desconocido.
        \vspace{0.1cm}
        \[ \mu \quad \sigma \quad p \quad N \]
      \end{block}
    \end{column}
    \begin{column}{0.48\textwidth}
      \begin{exampleblock}{Estadístico}
        \footnotesize
        Número que se calcula con la \azul{muestra}. Se conoce, pero cambia con
        cada muestra.
        \vspace{0.1cm}
        \[ \bar{x} \quad s \quad \hat{p} \quad n \]
      \end{exampleblock}
    \end{column}
  \end{columns}

  \vspace{0.3cm}
  \begin{alertblock}{Regla mnemotécnica}
    \footnotesize
    \destacar{Letras griegas para la población, letras latinas para la
    muestra.} Población y parámetro empiezan por P; muestra y muestral, por M.
  \end{alertblock}

  \vspace{0.1cm}
  {\footnotesize\color{uniGris} En este curso trabajamos casi siempre con
  estadísticos: describimos lo que tenemos a la mano.}
\end{frame}

\seccion{1.3 Tipos de variables}

\begin{frame}{Clasificación de las variables}
  \centering
  \begin{tikzpicture}[scale=0.9,every node/.style={font=\scriptsize,align=center},
                      nodo/.style={draw=uniAzul,fill=uniAzul!10,rounded corners,
                                   inner sep=3pt,minimum width=2.1cm}]
    \node[nodo] (v) at (0,0) {Variable};
    \node[nodo] (cl) at (-3,-1.3) {Cualitativa\\(categórica)};
    \node[nodo] (cn) at (3,-1.3) {Cuantitativa\\(numérica)};
    \node[nodo] (no) at (-4.6,-2.8) {Nominal};
    \node[nodo] (or) at (-1.6,-2.8) {Ordinal};
    \node[nodo] (di) at (1.5,-2.8) {Discreta};
    \node[nodo] (co) at (4.5,-2.8) {Continua};
    \draw[uniGris] (v)--(cl); \draw[uniGris] (v)--(cn);
    \draw[uniGris] (cl)--(no); \draw[uniGris] (cl)--(or);
    \draw[uniGris] (cn)--(di); \draw[uniGris] (cn)--(co);
  \end{tikzpicture}

  \vspace{0.3cm}
  \begin{alertblock}{La pregunta que decide}
    \footnotesize
    \destacar{¿El valor es un número con el que tiene sentido hacer
    operaciones?} Si sí, es cuantitativa. Si es una etiqueta, es cualitativa,
    aunque esté codificada con números (1 = hombre, 2 = mujer sigue siendo
    cualitativa).
  \end{alertblock}
\end{frame}

\begin{frame}{Los cuatro tipos, con ejemplos}
  \footnotesize
  \begin{center}
  \begin{tabular}{@{}p{2.0cm}p{4.6cm}p{4.6cm}@{}}
    \toprule
    \textbf{Tipo} & \textbf{Característica} & \textbf{Ejemplos} \\
    \midrule
    Nominal & Categorías sin orden natural &
      Canal de compra, ciudad, género, estado civil \\
    \addlinespace
    Ordinal & Categorías \azul{con} orden, pero sin distancia medible &
      Nivel de satisfacción, estrato, nivel educativo \\
    \addlinespace
    Discreta & Números que resultan de \azul{contar}; valores aislados &
      Número de hijos, pedidos por día, quejas \\
    \addlinespace
    Continua & Números que resultan de \azul{medir}; cualquier valor en un
      intervalo & Gasto, peso, tiempo, estatura \\
    \bottomrule
  \end{tabular}
  \end{center}

  \vspace{0.2cm}
  {\footnotesize\color{uniGris} Truco: si entre dos valores posibles siempre
  cabe otro, es continua. Entre 2 y 3 hijos no cabe nada; entre 2 y 3 kg cabe
  2{,}5.}
\end{frame}

\begin{frame}{Números que no son cantidades}
  \scriptsize
  \begin{center}
  \begin{tabular}{@{}p{3.4cm}p{3.2cm}p{2.4cm}p{2.8cm}@{}}
    \toprule
    \textbf{Variable} & \textbf{Se registra como} & \textbf{¿Es número?} &
    \textbf{Tipo real} \\
    \midrule
    Código de sucursal & 1, 2, 3 & Parece & Cualitativa nominal \\
    Género & 1 = hombre, 2 = mujer & Parece & Cualitativa nominal \\
    Estrato & 1 a 6 & Parece & Cualitativa ordinal \\
    Satisfacción & 1 a 5 & Parece & Cualitativa ordinal \\
    Número de hijos & 0, 1, 2, \dots & \azul{Sí} & Cuantitativa discreta \\
    Gasto mensual & 227{,}5 & \azul{Sí} & Cuantitativa continua \\
    \bottomrule
  \end{tabular}
  \end{center}

  \vspace{0.1cm}
  \begin{alertblock}{La prueba de los tres segundos}
    \scriptsize
    Pregúntese si \destacar{sumar dos valores produce algo con sentido}.
    ``Sucursal 1 + sucursal 2 = sucursal 3'' no significa nada: el número es
    una \azul{etiqueta}. En cambio ``2 hijos + 3 hijos = 5 hijos'' sí, y por
    eso una es cualitativa y la otra cuantitativa.
  \end{alertblock}
\end{frame}

\seccion{1.4 Escalas de medición}

\begin{frame}{Las cuatro escalas}
  \scriptsize
  \begin{center}
  \begin{tabular}{@{}p{1.8cm}p{3.2cm}p{2.9cm}p{2.9cm}@{}}
    \toprule
    \textbf{Escala} & \textbf{Permite} & \textbf{No permite} & \textbf{Ejemplo} \\
    \midrule
    Nominal & Clasificar, contar frecuencias & Ordenar ni operar &
      Canal de compra \\
    \addlinespace
    Ordinal & Clasificar y \azul{ordenar} & Restar: las distancias no son
      iguales & Satisfacción: malo, regular, bueno \\
    \addlinespace
    Intervalar & Ordenar y \azul{restar} & Dividir: el cero es arbitrario &
      Temperatura en $^\circ$C, año calendario \\
    \addlinespace
    De razón & Ordenar, restar y \azul{dividir} & --- (cero absoluto) &
      Gasto, peso, edad, tiempo \\
    \bottomrule
  \end{tabular}
  \end{center}

  \vspace{0.2cm}
  \begin{alertblock}{Por qué importa el cero absoluto}
    \footnotesize
    En razón, 40 mil es \destacar{el doble} de 20 mil. En intervalar no:
    $20^\circ$C no es el doble de calor que $10^\circ$C, porque el $0^\circ$C
    es una convención, no la ausencia de temperatura.
  \end{alertblock}
\end{frame}

\begin{frame}{Qué se puede calcular en cada escala}
  \footnotesize
  \begin{center}
  \begin{tabular}{@{}lcccc@{}}
    \toprule
    \textbf{Medida} & \textbf{Nominal} & \textbf{Ordinal} &
    \textbf{Intervalar} & \textbf{Razón} \\
    \midrule
    Frecuencias y moda      & Sí & Sí & Sí & Sí \\
    Mediana y percentiles   & No & Sí & Sí & Sí \\
    Media y desviación      & No & No & Sí & Sí \\
    Coeficiente de variación& No & No & No & Sí \\
    \bottomrule
  \end{tabular}
  \end{center}

  \vspace{0.25cm}
  \begin{alertblock}{Esta tabla es el mapa del curso}
    \footnotesize
    Cada vez que veamos una medida nueva, la pregunta será la misma:
    \destacar{¿la escala de esta variable me permite calcularla?} Promediar
    códigos de ciudad (1 = Bello, 2 = Medellín) da un número, pero no significa
    nada.
  \end{alertblock}
\end{frame}

\begin{frame}{El instrumento decide la escala}
  \scriptsize
  La misma característica cambia de escala según \destacar{cómo se pregunte}:
  \vspace{0.1cm}
  \begin{center}
  \begin{tabular}{@{}p{5.8cm}p{3.0cm}p{3.0cm}@{}}
    \toprule
    \textbf{Cómo se pregunta la edad} & \textbf{Tipo} & \textbf{Escala} \\
    \midrule
    ``¿Cuántos años cumplidos tiene?'' & Cuantitativa discreta & Razón \\
    ``¿En qué rango está: 18--25, 26--40, 41+?'' & Cualitativa ordinal &
      Ordinal \\
    ``¿Es usted mayor de edad? Sí / No'' & Cualitativa nominal & Nominal \\
    \bottomrule
  \end{tabular}
  \end{center}

  \vspace{0.15cm}
  \begin{alertblock}{Lo que se pierde en cada paso}
    \scriptsize
    De arriba abajo se \destacar{pierde información y nunca se recupera}: de la
    edad exacta se puede construir el rango, pero del rango no se puede
    recuperar la edad. \azul{Regla práctica: registre siempre en la escala más
    fina que pueda} y agrupe después, si hace falta.
  \end{alertblock}
\end{frame}

\begin{frame}{De la escala al gráfico y a la medida}
  \scriptsize
  \begin{center}
  \begin{tabular}{@{}p{2.0cm}p{3.2cm}p{3.2cm}p{3.4cm}@{}}
    \toprule
    \textbf{Escala} & \textbf{Tabla} & \textbf{Gráfico} &
    \textbf{Medida de centro} \\
    \midrule
    Nominal & Frecuencias, sin acumular &
      Barras, circular & Solo la moda \\
    \addlinespace
    Ordinal & Frecuencias \azul{con} acumuladas &
      Barras en orden natural & Mediana o moda \\
    \addlinespace
    Intervalar & Agrupada en clases &
      Histograma, polígono, ojiva & Media, mediana, moda \\
    \addlinespace
    De razón & Agrupada en clases &
      Histograma, polígono, ojiva & Todas, más el $CV$ \\
    \bottomrule
  \end{tabular}
  \end{center}

  \vspace{0.1cm}
  \begin{alertblock}{Este cuadro es el índice del curso}
    \scriptsize
    La columna \azul{Tabla} es la sesión 3; la de \azul{Gráfico}, la sesión 4;
    la de \azul{Medida}, toda la unidad 2. \destacar{Todo lo que sigue depende
    de haber clasificado bien la variable hoy.}
  \end{alertblock}
\end{frame}

\begin{frame}{Ejercicio 1 — en clase}
  \begin{alertblock}{En parejas (8 minutos)}
    Clasifique cada variable (cualitativa nominal / ordinal, cuantitativa
    discreta / continua) e indique su escala de medición.

    \vspace{0.25cm}
    \begin{columns}[T,onlytextwidth]
      \begin{column}{0.48\textwidth}
        \begin{enumerate}
          \item Gasto mensual del cliente
          \item Canal de compra
          \item Número de compras del mes
          \item Nivel de satisfacción (1 a 5)
        \end{enumerate}
      \end{column}
      \begin{column}{0.48\textwidth}
        \begin{enumerate}
          \setcounter{enumi}{4}
          \item Estrato socioeconómico
          \item Temperatura del local en $^\circ$C
          \item Tiempo de espera en minutos
          \item Código de la sucursal
        \end{enumerate}
      \end{column}
    \end{columns}
  \end{alertblock}
\end{frame}

\begin{frame}{Ejercicio 1 — solución}
  \scriptsize
  \begin{center}
  \begin{tabular}{@{}clll@{}}
    \toprule
    \textbf{\#} & \textbf{Variable} & \textbf{Tipo} & \textbf{Escala} \\
    \midrule
    1 & Gasto mensual        & Cuantitativa continua & Razón \\
    2 & Canal de compra      & Cualitativa nominal   & Nominal \\
    3 & Número de compras    & Cuantitativa discreta & Razón \\
    4 & Nivel de satisfacción& Cualitativa ordinal   & Ordinal \\
    5 & Estrato              & Cualitativa ordinal   & Ordinal \\
    6 & Temperatura ($^\circ$C) & Cuantitativa continua & Intervalar \\
    7 & Tiempo de espera     & Cuantitativa continua & Razón \\
    8 & Código de sucursal   & Cualitativa nominal   & Nominal \\
    \bottomrule
  \end{tabular}
  \end{center}

  \vspace{0.2cm}
  \begin{alertblock}{Las dos trampas}
    \footnotesize
    El \destacar{4 y el 5} vienen con números, pero la distancia entre
    satisfacción 1 y 2 no es la misma que entre 4 y 5: son ordinales. El
    \destacar{8} es un número que solo identifica: es nominal, y promediarlo no
    tiene sentido.
  \end{alertblock}
\end{frame}

\begin{frame}{Ejercicio 2 — en clase}
  \begin{alertblock}{Individual (7 minutos)}
    Para cada afirmación, diga si es \azul{parámetro} o \azul{estadístico} y
    escriba el símbolo correspondiente.

    \vspace{0.25cm}
    \begin{enumerate}
      \item El gasto promedio de los 40 clientes encuestados es 42{,}5.
      \item El gasto promedio de los 1\,200 clientes registrados.
      \item La desviación estándar calculada con la muestra.
      \item La proporción real de clientes que usan la app.
      \item ¿Cuál de los cuatro se conoce con certeza y por qué?
    \end{enumerate}
  \end{alertblock}
\end{frame}

\begin{frame}{Ejercicio 2 — solución}
  \begin{columns}[T,onlytextwidth]
    \begin{column}{0.52\textwidth}
      \begin{block}{Respuestas}
        \footnotesize
        \begin{enumerate}
          \item Estadístico: $\bar{x}=42{,}5$
          \item Parámetro: $\mu$
          \item Estadístico: $s$
          \item Parámetro: $p$
        \end{enumerate}
      \end{block}
    \end{column}
    \begin{column}{0.44\textwidth}
      \begin{exampleblock}{Punto 5}
        \footnotesize
        Se conocen \azul{1 y 3}, los estadísticos, porque se calculan con datos
        que sí tenemos.

        \vspace{0.15cm}
        Los parámetros exigirían medir a los 1\,200 clientes.
      \end{exampleblock}
    \end{column}
  \end{columns}

  \vspace{0.2cm}
  {\footnotesize\color{uniGris} En este curso describimos la muestra y nos
  detenemos ahí: calculamos $\bar{x}$, no $\mu$. \destacar{Los parámetros se
  nombran para saber qué no se está calculando.}}
\end{frame}

\begin{frame}{Ejercicio 3 — en clase}
  \begin{alertblock}{Individual (9 minutos)}
    Una encuesta de clima laboral recoge estas cuatro variables:

    \vspace{0.2cm}
    \centering\footnotesize
    \begin{tabular}{@{}ll@{}}
      \toprule
      \textbf{Variable} & \textbf{Cómo se registró} \\
      \midrule
      Área de trabajo        & Ventas / Operaciones / Administración \\
      Satisfacción general   & 1 = muy baja \dots\ 5 = muy alta \\
      Antigüedad             & Años cumplidos \\
      Salario mensual        & Pesos \\
      \bottomrule
    \end{tabular}

    \vspace{0.2cm}
    \raggedright
    \begin{enumerate}
      \item Clasifique cada una y diga su escala.
      \item ¿En cuáles tiene sentido la \emph{media}? ¿Y la \emph{moda}?
      \item Un informe reporta ``satisfacción promedio = 3{,}7''.
            ¿Es correcto? Justifique.
    \end{enumerate}
  \end{alertblock}
\end{frame}

\begin{frame}{Ejercicio 3 — solución}
  \scriptsize
  \begin{center}
  \begin{tabular}{@{}lllcc@{}}
    \toprule
    \textbf{Variable} & \textbf{Tipo} & \textbf{Escala} &
    \textbf{¿Media?} & \textbf{¿Moda?} \\
    \midrule
    Área de trabajo      & Cualitativa nominal   & Nominal & No & Sí \\
    Satisfacción (1--5)  & Cualitativa ordinal   & Ordinal & No & Sí \\
    Antigüedad           & Cuantitativa discreta & Razón   & Sí & Sí \\
    Salario mensual      & Cuantitativa continua & Razón   & Sí & Sí \\
    \bottomrule
  \end{tabular}
  \end{center}

  \vspace{0.2cm}
  \begin{alertblock}{Punto 3: la discusión que importa}
    \footnotesize
    Estrictamente \destacar{no es correcto}: en una escala ordinal la distancia
    entre 1 y 2 no es necesariamente igual a la que hay entre 4 y 5, así que
    sumar y dividir no está justificado.

    En la práctica se hace mucho, y no siempre es indefendible si la escala
    tiene 5 o más niveles. Lo \azul{seguro} es reportar la \azul{mediana} y la
    distribución de frecuencias.
  \end{alertblock}
\end{frame}

\begin{frame}{Ejercicio 4 — en clase (taller principal)}
  \begin{alertblock}{En grupos de 3 (20 minutos)}
    \scriptsize
    Una EPS arma una base con estas ocho variables de sus afiliados:

    \vspace{0.12cm}
    \centering
    \begin{tabular}{@{}llll@{}}
      \toprule
      1. Número de afiliación & 2. Edad en años & 3. Sexo &
      4. Estrato (1 a 6) \\
      5. Peso en kg & 6. Ciudad & 7. Consultas del año &
      8. Estado del caso \\
      \bottomrule
    \end{tabular}
    {\tiny (8: leve / moderado / grave)}

    \vspace{0.18cm}
    \raggedright
    \begin{enumerate}\setlength{\itemsep}{-0.06cm}
      \item Clasifique cada una: tipo y escala.
      \item ¿En cuáles tiene sentido la media? ¿En cuáles la mediana?
      \item ¿En cuáles se pueden usar frecuencias acumuladas?
      \item Un informe reporta ``estrato promedio 3,4''. ¿Es defendible?
      \item ¿Cuál de las ocho \destacar{no debería} entrar en ningún análisis
            estadístico, y por qué?
    \end{enumerate}
  \end{alertblock}
\end{frame}

\begin{frame}{Ejercicio 4 — solución}
  \tiny
  \begin{center}
  \begin{tabular}{@{}clllccc@{}}
    \toprule
    \textbf{\#} & \textbf{Variable} & \textbf{Tipo} & \textbf{Escala} &
    \textbf{¿Media?} & \textbf{¿Mediana?} & \textbf{¿Acumular?} \\
    \midrule
    1 & N.\ de afiliación & Cualitativa nominal & Nominal & No & No & No \\
    2 & Edad (años)       & Cuantitativa discreta & Razón & Sí & Sí & Sí \\
    3 & Sexo              & Cualitativa nominal & Nominal & No & No & No \\
    4 & Estrato           & Cualitativa ordinal & Ordinal & No & Sí & Sí \\
    5 & Peso (kg)         & Cuantitativa continua & Razón & Sí & Sí & Sí \\
    6 & Ciudad            & Cualitativa nominal & Nominal & No & No & No \\
    7 & Consultas del año & Cuantitativa discreta & Razón & Sí & Sí & Sí \\
    8 & Estado del caso   & Cualitativa ordinal & Ordinal & No & Sí & Sí \\
    \bottomrule
  \end{tabular}
  \end{center}

  \vspace{0.1cm}
  \begin{alertblock}{Puntos 4 y 5}
    \tiny
    \textbf{4.} Estrictamente \destacar{no}: el estrato es ordinal y la
    distancia entre 1 y 2 no es la misma que entre 5 y 6. Lo defendible es
    reportar la \azul{mediana} y la distribución de frecuencias.

    \vspace{0.06cm}
    \textbf{5.} El \destacar{número de afiliación}: identifica, no mide. Se usa
    para enlazar registros, nunca para calcular. Promediar cédulas es el
    ejemplo clásico de cifra correcta y sin significado.
  \end{alertblock}
\end{frame}

\begin{frame}{Errores frecuentes en esta sesión}
  \scriptsize
  \begin{center}
  \begin{tabular}{@{}p{5.4cm}p{6.4cm}@{}}
    \toprule
    \textbf{Error} & \textbf{Corrección} \\
    \midrule
    Tratar como cuantitativa toda variable con números &
      Códigos y escalas de 1 a 5 son \azul{etiquetas ordenadas} \\
    \addlinespace
    Promediar una variable ordinal sin advertirlo &
      Se reporta la mediana; si se usa la media, se justifica \\
    \addlinespace
    Confundir discreta con continua &
      Si entre dos valores siempre cabe otro, es continua \\
    \addlinespace
    Creer que ``intervalar'' y ``de razón'' son lo mismo &
      El cero de razón es \azul{ausencia}; el intervalar es una convención \\
    \addlinespace
    Ordenar alfabéticamente una variable ordinal &
      Se respeta el orden natural de la escala \\
    \addlinespace
    Usar $\bar{x}$ para la población &
      $\bar{x}$ y $s$ son de la muestra; $\mu$ y $\sigma$, de la población \\
    \bottomrule
  \end{tabular}
  \end{center}
\end{frame}

\begin{frame}[fragile]{Tipos y escalas, en R}
  \scriptsize
\begin{lstlisting}[style=Rmini]
gasto     <- c(112, 127, 138)
canal     <- c("Tienda", "App", "Web")
estrato   <- c(3, 2, 5)

class(gasto)      # "numeric": cuantitativa
class(canal)      # "character": R aun no sabe que es cualitativa

canal   <- factor(canal)                                  # nominal
estrato <- factor(estrato, levels = 1:6, ordered = TRUE)  # ORDINAL
str(estrato)
mean(as.numeric(estrato))   # R lo calcula... pero NO significa nada
\end{lstlisting}
\begin{lstlisting}[style=Rout]
[1] "numeric"
[1] "character"
 Ord.factor w/ 6 levels "1"<"2"<"3"<..: 3 2 5
[1] 3.333333
\end{lstlisting}

  \vspace{-0.12cm}
  \begin{alertblock}{La última línea es la lección de hoy}
    \tiny
    R \destacar{sí} calcula el promedio del estrato y devuelve 3,33 sin una
    sola advertencia. \azul{La herramienta no distingue escalas: usted sí.} Por
    eso se declara \fx{ordered = TRUE}: para que el orden quede escrito y
    quien lea el código sepa qué se puede y qué no se puede hacer.
  \end{alertblock}
\end{frame}

\begin{frame}{Cierre de la sesión}
  \begin{block}{Lo que debe quedar claro hoy}
    \begin{itemize}
      \item Parámetro describe la población ($\mu$, $\sigma$, $p$);
            estadístico describe la muestra ($\bar x$, $s$, $\hat p$).
      \item Cualitativa (nominal u ordinal) frente a cuantitativa (discreta o
            continua).
      \item Cuatro escalas: nominal, ordinal, intervalar y de razón.
      \item La escala decide qué tabla, qué gráfico y qué medida son válidos.
    \end{itemize}
  \end{block}

  \vspace{0.2cm}
  \begin{alertblock}{Trabajo independiente para la próxima sesión}
    \footnotesize
    Resolver los \destacar{20 ejercicios de la lámina siguiente}. En la sesión
    3 organizamos los datos en \azul{tablas de frecuencia}, y la escala de la
    variable será lo primero que preguntemos.
  \end{alertblock}
\end{frame}

\begin{frame}{Ejercicios propuestos}
  \vspace{-0.1cm}
  \begin{beamercolorbox}[wd=\textwidth,sep=0.10cm]{block body}
    \centering\scriptsize
    \textbf{1--12:} clasifique el tipo de variable y la escala. \quad
    \textbf{13--20:} parámetro o estadístico, y conceptos.
  \end{beamercolorbox}
  \listatiny
  \vspace{0.15cm}
  \begin{columns}[T,onlytextwidth]
    \begin{column}{0.47\textwidth}
      \begin{enumerate}\setlength{\itemsep}{0.16cm}
        \item Edad en años cumplidos
        \item Color favorito
        \item Número de hijos
        \item Salario mensual
        \item Nivel educativo alcanzado
        \item Número de cédula
        \item Estatura en metros
        \item Calificación de 1 a 5 estrellas
        \item Ciudad de residencia
        \item Año de fundación de la empresa
      \end{enumerate}
    \end{column}
    \begin{column}{0.50\textwidth}
      \begin{enumerate}\setlength{\itemsep}{0.16cm}
        \setcounter{enumi}{10}
        \item Quejas recibidas por día
        \item Tiempo de entrega en horas
        \item $\bar{x}$ calculado con 40 datos
        \item $\mu$ de los 1\,200 clientes
        \item $s$ de la muestra
        \item $\sigma$ de la población
        \item ¿Se puede calcular la media del estrato?
        \item ¿Se puede calcular la moda del canal de compra?
        \item ¿Se puede calcular la mediana del nivel de satisfacción?
        \item ¿Es $20^\circ$C el doble de $10^\circ$C?
      \end{enumerate}
    \end{column}
  \end{columns}
\end{frame}

\begin{frame}{Respuestas}{Autoevaluación de los ejercicios propuestos}
  \listatiny
  \vspace{0.15cm}
  \begin{columns}[T,onlytextwidth]
    \begin{column}{0.47\textwidth}
      \begin{enumerate}\setlength{\itemsep}{0.16cm}
        \item Cuantitativa discreta; razón
        \item Cualitativa nominal; nominal
        \item Cuantitativa discreta; razón
        \item Cuantitativa continua; razón
        \item Cualitativa ordinal; ordinal
        \item Cualitativa nominal; nominal
        \item Cuantitativa continua; razón
        \item Cualitativa ordinal; ordinal
        \item Cualitativa nominal; nominal
        \item Cuantitativa discreta; intervalar
      \end{enumerate}
    \end{column}
    \begin{column}{0.50\textwidth}
      \begin{enumerate}\setlength{\itemsep}{0.16cm}
        \setcounter{enumi}{10}
        \item Cuantitativa discreta; razón
        \item Cuantitativa continua; razón
        \item Estadístico
        \item Parámetro
        \item Estadístico
        \item Parámetro
        \item No: es ordinal
        \item Sí: la moda vale en nominal
        \item Sí: la mediana vale en ordinal
        \item No: el cero de $^\circ$C es arbitrario
      \end{enumerate}
    \end{column}
  \end{columns}
\end{frame}

\begin{frame}{Referencias}
  \footnotesize
  \begin{itemize}
    \item[{[1]}] Monroy Saldívar, S. (2008). \textit{Estadística descriptiva}.
      Instituto Politécnico Nacional.
      \texttt{https://elibro.net/es/lc/uniminuto/titulos/74722}
    \item[{[2]}] Ordóñez Fernández, F. F. \& González Fernández, J. (2021).
      \textit{Estadística descriptiva paso a paso}. Ediciones Pirámide.
      \texttt{https://elibro.net/es/lc/uniminuto/titulos/215449}
    \item[{[3]}] Llinás Solano, H. (2014). \textit{Introducción a la
      estadística con aplicaciones en ciencias sociales}. Universidad del
      Norte. \texttt{https://elibro.net/es/lc/uniminuto/titulos/70062}
  \end{itemize}

  \vspace{0.35cm}
  \begin{center}
    {\Large\color{uniAzul}\textbf{¿Preguntas?}}
  \end{center}
\end{frame}

\end{document}
